\documentclass{article}
\usepackage[utf8]{inputenc}
\usepackage[margin=1in]{geometry}
\usepackage{xcolor}
\usepackage{listings}

\definecolor{codegreen}{rgb}{0,0.6,0}
\definecolor{codegray}{rgb}{0.5,0.5,0.5}
\definecolor{codepurple}{rgb}{0.58,0,0.82}
\definecolor{backcolour}{rgb}{0.95,0.95,0.92}

\lstdefinestyle{sqlstyle}{
    backgroundcolor=\color{backcolour},
    commentstyle=\color{codegreen},
    keywordstyle=\color{blue},
    numberstyle=\tiny\color{codegray},
    stringstyle=\color{codepurple},
    basicstyle=\ttfamily\footnotesize,
    breakatwhitespace=false,
    breaklines=true,
    captionpos=b,
    keepspaces=true,
    numbers=left,
    numbersep=5pt,
    showspaces=false,
    showstringspaces=false,
    showtabs=false,
    tabsize=2
}
\lstset{style=sqlstyle}

\title{Lab 7: Roommate Expense Splitter}
\author{Name: Rahul Paul \\ Roll No: 202411078 \\ DBMS Lab-7}
\date{\today}

\begin{document}

\maketitle

\section{Table Creation}
\begin{lstlisting}[language=SQL]
DROP TABLE IF EXISTS expenses;
DROP TABLE IF EXISTS roommates;

CREATE TABLE roommates (
    id INT PRIMARY KEY,
    name VARCHAR(50)
);

CREATE TABLE expenses (
    id INT PRIMARY KEY,
    payer_id INT REFERENCES roommates(id),
    amount INT,
    description VARCHAR(100)
);
\end{lstlisting}

\section{Sample Data}
\begin{lstlisting}[language=SQL]
INSERT INTO roommates (id, name) VALUES
(1, 'Aarav'),
(2, 'Diya'),
(3, 'Kabir'),
(4, 'Meera');

INSERT INTO expenses (id, payer_id, amount, description) VALUES
(1, 1, 1800, 'Rent and Groceries'),
(2, 2, 800, 'Electricity'),
(3, 3, 1000, 'WiFi and Dinner'),
(4, 4, 400, 'Cleaning Supplies');
\end{lstlisting}

\section{Queries and Outputs}

\subsection*{Query A: Show each expense with roommate name}
\begin{lstlisting}[language=SQL]
SELECT r.name, e.description, e.amount
FROM roommates r
JOIN expenses e ON r.id = e.payer_id;
\end{lstlisting}
\textbf{Output A:} Paste output here.

\subsection*{Query B: Show total spent by each roommate}
\begin{lstlisting}[language=SQL]
SELECT r.name, SUM(e.amount) AS total_spent
FROM roommates r
JOIN expenses e ON r.id = e.payer_id
GROUP BY r.name;
\end{lstlisting}
\textbf{Output B:} Paste output here.

\subsection*{Query C: Show total expense pool and fair share per person}
\begin{lstlisting}[language=SQL]
SELECT 
    SUM(amount) AS total_pool, 
    SUM(amount) / (SELECT COUNT(*) FROM roommates) AS fair_share 
FROM expenses;
\end{lstlisting}
\textbf{Output C:} Paste output here.

\subsection*{Query D: Show roommates who still need to pay}
\begin{lstlisting}[language=SQL]
SELECT 
    r.name, 
    SUM(e.amount) AS amount_paid,
    (SELECT SUM(amount)/COUNT(*) FROM roommates) - SUM(e.amount) AS amount_to_pay
FROM roommates r
JOIN expenses e ON r.id = e.payer_id
GROUP BY r.name
HAVING SUM(e.amount) < (SELECT SUM(amount)/COUNT(*) FROM roommates);
\end{lstlisting}
\textbf{Output D:} Paste output here.

\subsection*{Query E: Show expenses greater than 500}
\begin{lstlisting}[language=SQL]
SELECT r.name, e.description, e.amount
FROM roommates r
JOIN expenses e ON r.id = e.payer_id
WHERE e.amount > 500;
\end{lstlisting}
\textbf{Output E:} Paste output here.

\subsection*{Query F: Show expenses related to groceries}
\begin{lstlisting}[language=SQL]
SELECT r.name, e.description, e.amount
FROM roommates r
JOIN expenses e ON r.id = e.payer_id
WHERE e.description LIKE '%Groceries%';
\end{lstlisting}
\textbf{Output F:} Paste output here.

\subsection*{Query G: Count expense entries per roommate}
\begin{lstlisting}[language=SQL]
SELECT r.name, COUNT(e.id) AS number_of_purchases
FROM roommates r
JOIN expenses e ON r.id = e.payer_id
GROUP BY r.name;
\end{lstlisting}
\textbf{Output G:} Paste output here.

\subsection*{Query H: Show the highest single expense}
\begin{lstlisting}[language=SQL]
SELECT r.name, e.description, e.amount
FROM roommates r
JOIN expenses e ON r.id = e.payer_id
ORDER BY e.amount DESC
LIMIT 1;
\end{lstlisting}
\textbf{Output H:} Paste output here.

\end{document}
